%------------------------
% Resume Template
% Author : Anubhav Singh
% Github : https://github.com/xprilion
% License : MIT
%------------------------

\documentclass[a4paper]{article}

\RequirePackage{fix-cm}
\usepackage{latexsym}
\usepackage[empty]{fullpage}
\usepackage{titlesec}
\usepackage{marvosym}
\usepackage[usenames,dvipsnames]{color}
\usepackage{verbatim}
\usepackage{enumitem}
\usepackage[pdftex,hidelinks]{hyperref}
\usepackage{fancyhdr}

\pagestyle{fancy}
\fancyhf{}
\fancyfoot{}
\renewcommand{\headrulewidth}{0pt}
\renewcommand{\footrulewidth}{0pt}

\usepackage[T1]{fontenc}
\usepackage{lmodern}
\input{glyphtounicode}
\pdfgentounicode=1
\usepackage{microtype}
\DisableLigatures[f]{encoding = *, family = *}

% Margins
\addtolength{\oddsidemargin}{-0.530in}
\addtolength{\evensidemargin}{-0.375in}
\addtolength{\textwidth}{1in}
\addtolength{\topmargin}{-0.7in}
\addtolength{\textheight}{1.35in}

\urlstyle{rm}
\raggedbottom
\raggedright
\setlength{\tabcolsep}{0in}

% Section headings
\titleformat{\section}{\vspace{-8pt}\scshape\raggedright\large}{}{0em}{}[\color{black}\titlerule\vspace{-5pt}]

%-------------------------
% Custom commands
\newcommand{\resumeItem}[2]{\item\fontsize{9.5}{11.3}\selectfont{\textbf{#1}{: #2}}}
\newcommand{\bul}[1]{\item\fontsize{9.5}{11.3}\selectfont{#1}}
\newcommand{\resumeSubheading}[4]{
  \vspace{1pt}\item
    \begin{tabular*}{0.985\textwidth}{l@{\extracolsep{\fill}}r}
      \textbf{#1} & #2 \\
      \textit{#3} & \textit{#4} \\
    \end{tabular*}\vspace{-1pt}
}
\newcommand{\resumeSubItem}[2]{\resumeItem{#1}{#2}\vspace{-1pt}}
\newcommand{\resumeSubHeadingListStart}{\begin{itemize}[leftmargin=0pt,label={},itemsep=1pt,topsep=2pt]}
\newcommand{\resumeSubHeadingListEnd}{\end{itemize}}
\newcommand{\resumeItemListStart}{\begin{itemize}[leftmargin=12pt,itemsep=0.6pt,topsep=1pt,parsep=0pt]}
\newcommand{\resumeItemListEnd}{\end{itemize}\vspace{-3pt}}

%-----------------------------
\begin{document}
\fontfamily{cmss}\selectfont

%---------- HEADING ----------
\begin{tabular*}{\textwidth}{l@{\extracolsep{\fill}}r}
  \textbf{{\LARGE Rahul Muralidhara Koushik}} & San Francisco Bay Area \\
  Email: \href{mailto:rkoushik@asu.edu}{rkoushik@asu.edu} & Mobile:~+1-480-791-6597 \\
  \href{https://github.com/Fluff18}{GitHub: github.com/Fluff18} \;|\; \href{https://fluff18.github.io/}{Portfolio: fluff18.github.io} & LinkedIn: \href{https://www.linkedin.com/in/rahulkoushik/}{linkedin.com/in/rahulkoushik} \\
\end{tabular*}

\vspace{2pt}
{\fontsize{9.5}{11.3}\selectfont \textbf{Software Engineer \textbar{} Distributed Systems, Backend and AI Infrastructure.} Three years building Go and Kubernetes services, observability and enterprise data migration for HPE's storage platform, now building production AI agent infrastructure on the same systems.}

%---------- EXPERIENCE ----------
\section{Experience}
\resumeSubHeadingListStart
  \resumeSubheading{Hewlett Packard Enterprise, Storage AI Team}{San Jose, US}{AI Engineer}{Feb 2026 -- Present}
  \resumeItemListStart
    \bul{Architected the multi-agent LLM harness every team agent runs on, with an MCP tool server, guardrails and tracing; it serves \textbf{6,300+ tool calls a week} across HPE Storage.}
    \bul{Built the memory layer for persistent agents: embeddings, vector search and a context graph linking \textbf{13,005} past bugs and \textbf{5,600} analyses by component, symptom and fix.}
    \bul{Built the agent evaluation platform in four weeks (\textbf{three production deployments}): rubric scoring against expert feedback, LLM-as-a-judge and tracing, now the release gate.}
  \resumeItemListEnd

  \resumeSubheading{Hewlett Packard Enterprise, Storage}{San Jose, US}{Cloud Software Engineer}{Jun 2023 -- Feb 2026}
  \resumeItemListStart
    \bul{Built cross-stack observability for a platform serving \textbf{10,000+ customers}, cutting typical issue resolution from \textbf{30 to 3 minutes}; failure-trend analysis raised first-time success from \textbf{55\% to 85\%}.}
    \bul{Implemented six enterprise data-migration features in Go across four storage array generations, over Fibre Channel and iSCSI.}
    \bul{Built a documentation agent that updates docs on every pull request, saving about \textbf{5 engineer-hours a week} per team.}
  \resumeItemListEnd

  \resumeSubheading{The Luminosity Lab, Arizona State University}{Tempe, US}{Research Aide}{Feb 2022 -- May 2023}
  \resumeItemListStart
    \bul{Led the iOS team on the \href{https://thetouchpointsolution.com/}{TouchPoint} apps; a telehealth flow saved doctors about 15 minutes per visit.}
  \resumeItemListEnd

  \resumeSubheading{Aruba, Hewlett Packard Enterprise}{Tempe, US}{Software Engineering Intern}{May 2022 -- Aug 2022}
  \resumeItemListStart
    \bul{Built a test-automation tool that configures network topologies, saving an estimated 150k engineering hours a year.}
  \resumeItemListEnd

  \resumeSubheading{Webex by Cisco}{Bengaluru, India}{Technical Undergraduate Intern}{Jan 2021 -- Jul 2021}
  \resumeItemListStart
    \bul{Built call and meeting scheduling APIs for the \href{https://developer.webex.com/docs/sdks/ios}{Webex iOS SDK v3} and the \href{https://github.com/webex/webex-ios-sdk-example}{Kitchen Sink} sample app.}
  \resumeItemListEnd
\resumeSubHeadingListEnd

%---------- SKILLS ----------
\section{Skills}
\resumeSubHeadingListStart
  \resumeSubItem{\textbullet\ Backend and languages}{Go, Python, TypeScript, gRPC, REST, FastAPI, SQL}
  \resumeSubItem{\textbullet\ Infrastructure and data}{Kubernetes, Helm, Docker, AWS, Kafka, Spark, PostgreSQL, Prometheus, Grafana, Linux, CI/CD}
  \resumeSubItem{\textbullet\ Storage and reliability}{Enterprise data migration, Fibre Channel and iSCSI, observability, log pipelines, debugging at release scale}
  \resumeSubItem{\textbullet\ AI systems}{LLM agents, MCP, RAG, evaluations, vLLM}
\resumeSubHeadingListEnd

%---------- PROJECTS ----------
\section{Selected Projects}
\resumeSubHeadingListStart
  \resumeSubItem{\textbullet\ Zuo, personal agent infrastructure (2026)}{Job queue, budgets and agent teams dispatched from my phone to cloud endpoints, a vLLM GPU and a llama.cpp laptop.}
  \resumeSubItem{\textbullet\ Catan Arena, agent evaluation testbed (2026)}{A learning agent graduates only when its win rate's 95\% Wilson interval clears the 25\% null.}
  \resumeSubItem{\textbullet\ Zzz, iOS app (2026)}{Gamified goal tracker released on the App Store, built solo in SwiftUI with Supabase.}
\resumeSubHeadingListEnd

%---------- HONORS ----------
\section{Honors}
\resumeSubHeadingListStart
  \bul{Three HPE hackathon wins (2024, 2026); two projects selected for HPE Fellow Review and BU-wide implementation \textbullet{} HPE Early Career Technologist (2024) \textbullet{} Patent application in progress (2026)}
\resumeSubHeadingListEnd

%---------- EDUCATION AND PUBLICATIONS ----------
\section{Education and Publications}
\resumeSubHeadingListStart
  \bul{\textbf{M.S. Computer Science}, Arizona State University, Tempe \hfill Aug 2021 -- May 2023}
  \bul{\textbf{B.Tech Computer Science and Engineering}, PES University, Bengaluru \hfill Aug 2017 -- Jul 2021}
  \bul{\textbf{Publications:} ViStA (IEEE ISPASS 2021); true random numbers from portable-computer entropy (IEEE CONECCT 2021); \href{https://doi.org/10.1007/978-981-33-4543-0_81}{PipeGAN} (Springer LNNS 171, 2021).}
\resumeSubHeadingListEnd

\end{document}
